\clearpage
\begin{center}
{\large SGA 7 --- 第 XVII 讲}\par
\vspace{2em}
{\Large\underline{Lefschetz 线束: 存在定理}}\par
\vspace{1.5em}
{\large N. M. Katz}
\end{center}

\section*{0. 引言}

本讲(在口头讨论班中先于讨论 Picard--Lefschetz 公式的第 XIV 和 XV 讲)介绍光滑射影概形的超平面截面的“Lefschetz 线束”, 并给出它们存在的条件. 下一讲将基于爆破簇上同调的已知结构, 对这些束作初等的上同调研究. 我们将引入的 Lefschetz 线束的某些假设经常成立, 这一事实将由第 XV 讲的 Picard--Lefschetz 公式得出; 此公式在性质上不如本讲及下一讲的结果初等.

本讲及下一讲的一些论述将在别处置于更一般的背景下展开(一般的底概形、用拟射影性假设代替射影性等): 特别是超平面截面束尤其是 Lefschetz 线束的上同调前研究见 EGA V, 而爆破概形的上同调研究见 Jouanolou 在 SGA 5 中的第 VII 讲. 读者还可在 EGA V 中找到以内蕴的图追踪取代本讲作者所偏爱的坐标方法的证明.

我要感谢 P. Deligne 和 P. Griffiths, 他们为本讲的构思提供了帮助. 除非明确另有说明, 全讲中 $k$ 均表示代数闭域.

\section*{1. 二次奇点回顾(参见 VI, 6)}

\subsubsection*{1.1.}
设 $Y$ 为纯维数 $n-1$ 的 $k$-概形.
\footnote{印本先写“维数为 $n$”, 而在下面的超曲面表示中又规定 $n=1+\dim Y$. 所列 Hessian 及正规形使用的是后一个约定.}
闭点 $y\in Y$ 称为 $Y$ 的\underline{普通二次奇点}, 如果:
\begin{enumerate}
\item[(a)] $Y$ 在 $y$ 处局部环的完备化 $\widehat{\mathcal O}_{Y,y}$ 同构于
\[
k[[T_1,\ldots,T_n]]/(Q(T)),
\]
即 $n=1+\dim Y$ 元形式幂级数环模由元素 $Q(T)$ 生成的理想的商;
\item[(b)] $Q$ 的最低次为二, 且当 $n\geq2$ 时, 由 $Q(T)$ 的二次齐次部分的消失所定义的 $\mathbb P^{n-1}_k$ 的子簇在 $k$ 上光滑; 当 $n=1$ 时, $Q(T)$ 的二次部分非零.
\end{enumerate}

若 $k$ 的特征不为二, \underline{或} $n$ 为偶数, 则条件 (b) 等价于:
\begin{enumerate}
\item[(b$'$)] $Q$ 的最低次为二, 且其 Hessian 矩阵
\[
\left(
  \frac{\partial^2Q}{\partial T_i\partial T_j}(0)
\right)_{1\leq i,j\leq n}
\]
可逆.
\end{enumerate}
这由光滑性的 Jacobian 判别法得出.

满足 (b) 的级数 $Q$ 的二次齐次部分 $Q_2$ 经过线性换元后可化为标准形
\begin{equation}\tag{1.1.1}
\left\{
\begin{aligned}
Q_2(T)&=\sum_{i=1}^{n/2}T_{2i-1}T_{2i}
  &&\text{若 $n$ 为偶数},\\[1ex]
Q_2(T)&=T_n^2+\sum_{i=1}^{(n-1)/2}T_{2i-1}T_{2i}
  &&\text{若 $n$ 为奇数}.
\end{aligned}
\right.
\end{equation}
这由二次型理论得出(Bourbaki, \emph{Alg\`ebre}, 第 X 章).

\subsubsection*{1.2.}
% TODO(核对): 英译及法文印本第 214 页均写 (a), (b); 非退化条件似应引用 (b'). 此处忠实保留底本, 不静默改正.
满足 (a) 和 (b) 的闭点 $y\in Y$ 称为\underline{非退化二次奇点}. 由于在特征二下 $n$ 为奇数时 Hessian 矩阵的行列式恒为零, 以 $n$ 个变量表示的非退化二次奇点仅当 $n$ 为偶数或 $k$ 的特征不为二时才存在; 在这些情形中, 它们恰为普通二次奇点.